\chapter{The plant model}
\label{ch:plant}

\code{pyplant/core/hybrid_tank.py} holds the physics as pure functions: SI
units throughout, no sockets, no sleeping, no global state. This chapter
derives every equation it implements.

\section{State, inputs and outputs}

The state vector is
\[
  \bm{x} = \begin{bmatrix} h_1 & h_2 & z \end{bmatrix}^\top,
\]
the two levels in metres and the valve stem position $z\in[0,1]$. The inputs
are the valve command $u\in[0,1]$ and the pump run command. The outputs are
the instrument readings of Table~\ref{tab:meas}.

\begin{table}[htbp]
\centering
\begin{tabular}{@{}llll@{}}
\toprule
Tag & Quantity & Formula & Span \\
\midrule
\tg{LT101} & tank 1 level & $1000\,h_1$ & 0--600\,mm \\
\tg{LT102} & tank 2 level & $1000\,h_2$ & 0--600\,mm \\
\tg{FT101} & pump discharge flow & $60\,000\,q_\text{in}$ & 0--20\,L/min \\
\tg{FT102} & tank 2 outlet flow & $60\,000\,q_\text{out}$ & 0--20\,L/min \\
\tg{PT101} & pump discharge pressure & $p_\text{pump}/1000$ & 0--100\,kPa \\
\tg{ZT101} & valve stem position & $100\,z$ & 0--100\,\% \\
\tg{PMP101_RUN} & pump running feedback & pump command & bool \\
\bottomrule
\end{tabular}
\caption{Measurements. \code{measure()} is the only function that converts SI
to engineering units.}
\label{tab:meas}
\end{table}

\section{Mass balances}

Each tank is a vertical cylinder of cross-section $A$. Conservation of volume
for an incompressible liquid gives
\begin{align}
  A_1 \frac{dh_1}{dt} &= q_\text{in} - q_{12} - q_\text{leak}, \label{eq:tank1}\\
  A_2 \frac{dh_2}{dt} &= q_{12} - q_\text{out}. \label{eq:tank2}
\end{align}
Both tanks are bounded: a level at the overflow lip $h_\text{max}=0.6$\,m
cannot rise further, and an empty tank cannot drain.

\section{Gravity drains: Torricelli's law}

Flow through an orifice under a head $h$ follows from Bernoulli's equation
between the free surface and the jet, $v=\sqrt{2gh}$, multiplied by an
effective discharge area $C_dA$ that lumps the geometric area and the
contraction coefficient:
\begin{equation}
  q = C_dA\,\sqrt{2gh}.
\end{equation}
This gives the interconnect, outlet and leak flows:
\begin{align}
  q_{12} &= C_dA_{12}\,g_{12}\,\sqrt{2gh_1}, &
  q_\text{out} &= C_dA_{2,\text{out}}\,\sqrt{2gh_2}, &
  q_\text{leak} &= C_dA_\text{leak}\,\sqrt{2gh_1},
\end{align}
where $g_{12}\le1$ is the fouling factor and $C_dA_\text{leak}\ge0$ the leak
area. Both are healthy at $g_{12}=1$, $C_dA_\text{leak}=0$.

\section{Pump and valve: an algebraic loop solved in closed form}

A centrifugal pump develops less pressure as it delivers more flow. A
quadratic curve is the standard approximation:
\begin{equation}
  p_\text{pump} = \Delta p_0\,\gamma\left(1-\left(\frac{q}{q_\text{max}}\right)^2\right),
\end{equation}
with shut-off pressure $\Delta p_0 = 60$\,kPa, run-out flow
$q_\text{max} = 25$\,L/min and wear factor $\gamma\le1$.

The control valve passes a flow proportional to the square root of the
pressure drop across it. The drop is the pump pressure minus the static head
of tank~1:
\begin{equation}
  q = k\sqrt{p_\text{pump}-\rho g h_1}, \qquad k = C_v\,\phi(z).
\end{equation}

The valve has an \emph{equal-percentage} inherent characteristic with
rangeability $R=30$: equal increments of stem travel change the flow by equal
percentages,
\begin{equation}
  \phi(z) = R^{\,z-1}, \qquad \phi(1) = 1,\ \phi(0) = 1/R,
\end{equation}
and below a seat position of 0.5\% the valve shuts tight. Under a falling
pump curve, an equal-percentage valve gives a roughly linear
\emph{installed} characteristic, which is why it is chosen.

Pump pressure depends on flow, and flow depends on pump pressure. Substituting
one into the other and solving the quadratic gives
\begin{equation}
  \boxed{\;q_\text{in} = k\sqrt{\dfrac{\Delta p_0\gamma - \rho g h_1}
        {1 + k^2\,\Delta p_0\gamma/q_\text{max}^2}}\;}
  \label{eq:qin}
\end{equation}
so no iteration is needed and the system is not stiff at a 50\,ms step. If the
valve is shut the pump dead-heads at $\Delta p_0\gamma$; if the pump is off,
$q_\text{in}=0$.

\begin{figure}[htbp]
\centering
\begin{subfigure}{0.4\textwidth}
\begin{tikzpicture}
\begin{axis}[width=\textwidth,height=3.8cm,xlabel={stem position $z$},
  ylabel={$\phi(z)$},grid=major,grid style={linegrey},
  xmin=0,xmax=1,ymin=0,ymax=1]
\addplot[accent,very thick,domain=0.005:1,samples=80]{30^(x-1)};
\addplot[slate,dashed,domain=0:1]{x};
\legend{equal \%, linear}
\end{axis}
\end{tikzpicture}
\caption{Inherent valve characteristic, $R=30$.}
\end{subfigure}\hfill
\begin{subfigure}{0.4\textwidth}
\begin{tikzpicture}
\begin{axis}[width=\textwidth,height=3.8cm,xlabel={flow (L/min)},
  ylabel={pressure (kPa)},grid=major,grid style={linegrey},
  xmin=0,xmax=25,ymin=0,ymax=65]
\addplot[accent,very thick,domain=0:25,samples=60]{60*(1-(x/25)^2)};
\addplot[badred,very thick,domain=0:25,samples=60]{0.85*60*(1-(x/25)^2)};
\addplot[only marks,mark=*,mark size=2pt] coordinates {(10,50.4)};
\legend{healthy,15\% wear,design point}
\end{axis}
\end{tikzpicture}
\caption{Pump curve, healthy and worn.}
\end{subfigure}
\caption{Component characteristics used by the plant model.}
\label{fig:curves}
\end{figure}

\section{Valve actuator}

The stem follows the command through a first-order lag with
$\tau_v = 2$\,s:
\begin{equation}
  \tau_v \frac{dz}{dt} = u - z .
\end{equation}

\section{Numerical integration}

The model is integrated with the classical fourth-order Runge--Kutta method at
$\Delta t = 50$\,ms:
\begin{align*}
  \bm k_1 &= f(\bm x_n,u), &
  \bm k_2 &= f(\bm x_n + \tfrac{\Delta t}{2}\bm k_1,u),\\
  \bm k_3 &= f(\bm x_n + \tfrac{\Delta t}{2}\bm k_2,u), &
  \bm k_4 &= f(\bm x_n + \Delta t\,\bm k_3,u),\\
  \bm x_{n+1} &= \bm x_n + \tfrac{\Delta t}{6}(\bm k_1+2\bm k_2+2\bm k_3+\bm k_4),
\end{align*}
after which levels are clipped to $[0,h_\text{max}]$ and $z$ to $[0,1]$. The
input is held constant within a step.

\section{Steady state}

At equilibrium $dh_1/dt = 0$. For a fixed valve position, inflow falls and
outflow rises monotonically with $h_1$, so the residual
$r(h_1) = q_\text{in} - q_{12} - q_\text{leak}$ has a single root, found by
bisection. Then $dh_2/dt=0$ gives tank 2 directly:
\begin{equation}
  h_2 = \frac{1}{2g}\left(\frac{q_{12}}{C_dA_{2,\text{out}}}\right)^2 .
\end{equation}
\code{steady_state()} is used to start the plant lined out and by the tests.

\section{Sizing and the design point}

\begin{table}[htbp]
\centering
\begin{tabular}{@{}lll@{}}
\toprule
Parameter (\code{Params} field) & Value & Meaning \\
\midrule
\code{area_tank1}, \code{area_tank2} & 0.020\,m$^2$ & tank cross-sections ($\approx$160\,mm diameter) \\
\code{height_max} & 0.600\,m & overflow lip \\
\code{pump_dp0} & 60\,000\,Pa & shut-off pressure \\
\code{pump_qmax} & $4.17\times10^{-4}$\,m$^3$/s & run-out flow, 25\,L/min \\
\code{valve_cv} & $2.429\times10^{-6}$ & valve coefficient fully open, m$^3$/(s$\sqrt{\text{Pa}}$) \\
\code{valve_rangeability} & 30 & equal-percentage $R$ \\
\code{valve_seat} & 0.005 & tight shut-off below this position \\
\code{tau_valve} & 2.0\,s & actuator lag \\
\code{cda_12} & $6.90\times10^{-5}$\,m$^2$ & tank 1 to tank 2 orifice \\
\code{cda_2out} & $9.00\times10^{-5}$\,m$^2$ & tank 2 drain \\
\code{pump_gain} & 1.0 & fault knob: $<1$ is impeller wear \\
\code{cda_12_gain} & 1.0 & fault knob: $<1$ is fouling \\
\code{cda_leak} & 0.0 & fault knob: leak area \\
\bottomrule
\end{tabular}
\caption{Plant parameters and their defaults.}
\label{tab:params}
\end{table}

At the design point the valve is 66\% open, the level is 300\,mm, the flow
10\,L/min and the pump pressure 50\,kPa. The usable level range is roughly 30
to 580\,mm, corresponding to valve positions from 0.30 to 0.79.

\section{Linearisation and the process time constant}

Tuning needs a linear model. Around an operating point $(\bar h_1,\bar z)$,
write $h_1 = \bar h_1 + \delta h$, $z = \bar z + \delta z$ and expand
equation~\eqref{eq:tank1} to first order:
\begin{equation}
  A_1 \frac{d\,\delta h}{dt}
   = \underbrace{\left(\frac{\partial q_\text{in}}{\partial h_1}
     - \frac{\partial q_{12}}{\partial h_1}\right)}_{-1/R_h}\delta h
   + \frac{\partial q_\text{in}}{\partial z}\,\delta z .
\end{equation}
Inflow falls as the back-pressure rises, so $\partial q_\text{in}/\partial h_1<0$,
and outflow rises with head, $\partial q_{12}/\partial h_1 = q_{12}/(2h_1)>0$;
the bracket is therefore negative and the tank is self-regulating. This is a
first-order system
\begin{equation}
  \tau \frac{d\,\delta h}{dt} + \delta h = K\,\delta z,
  \qquad \tau = A_1 R_h,\quad K = R_h\frac{\partial q_\text{in}}{\partial z}.
\end{equation}
An open-loop step test on the model at the design point gives
$K \approx 15.4$\,mm per \% of valve travel and $\tau\approx 72$\,s, the
values the controller is tuned for. The valve lag (2\,s) is small compared
with $\tau$ and is neglected in tuning.

\begin{figure}[htbp]
\centering
\scalebox{0.8}{%
\begin{tikzpicture}[node distance=12mm]
  \node[box] (u) {valve command\\$u$};
  \node[box,right=of u] (v) {actuator\\$\dfrac{1}{\tau_v s+1}$};
  \node[box,right=of v] (q) {pump + valve\\eq.~\eqref{eq:qin}};
  \node[hot,right=of q] (t1) {tank 1\\$A_1\dot h_1$};
  \node[box,right=of t1] (t2) {tank 2\\$A_2\dot h_2$};
  \draw[flow] (u) -- (v);
  \draw[flow] (v) -- node[lbl,above]{$z$} (q);
  \draw[flow] (q) -- node[lbl,above]{$q_\text{in}$} (t1);
  \draw[flow] (t1) -- node[lbl,above]{$q_{12}$} (t2);
  \draw[flow] (t1.south) -- ++(0,-8mm) -| node[lbl,pos=0.25,below]{$\rho g h_1$ back-pressure} (q.south);
  \draw[flow] (t2.east) -- ++(8mm,0) node[right]{$q_\text{out}$};
  \draw[flow,badred] (t1.north) -- ++(0,7mm) node[above]{$q_\text{leak}$};
\end{tikzpicture}}
\caption{Signal flow inside the plant model.}
\label{fig:plant-signal}
\end{figure}

\section{The simulation stepper}

\code{pyplant/runners/base.py} holds \code{Simulation}, the one loop body both
runners drive. It has three deliberately different rates:

\begin{center}
\begin{tabular}{@{}lll@{}}
\toprule
Rate & Default & What it governs \\
\midrule
\code{dt} & 50\,ms & integration accuracy \\
\code{control_period} & 100\,ms & how often the controller sees the plant \\
\code{log_period} & 1\,s & what the historian and datasets store \\
\bottomrule
\end{tabular}
\end{center}

The controller period must be a whole number of integration steps, and time
is computed from an integer step count rather than by repeated addition, so it
never accumulates rounding error. Between controller executions the inputs
are held: a genuine zero-order hold, as on a real PLC. Faults are applied in
\code{Simulation.integrate()}, the single place both runners share
(Chapter~\ref{ch:faults}).
